\chapter{PHÂN TÍCH YÊU CẦU VÀ DỮ LIỆU}
\label{chap:3-requirements}

% TODO: đoạn dẫn nhập chương

\section{Lĩnh vực ứng dụng: giám sát nông nghiệp}
\label{sec:3-domain}
% TODO

\section{Các bên liên quan}
\label{sec:3-stakeholders}
% TODO: quản lý trang trại, vận hành viên đồng ruộng, nhà phân tích dữ liệu, ...

\section{Yêu cầu nghiệp vụ}
\label{sec:3-br}

\begin{table}[H]
\centering
\caption{Danh sách yêu cầu nghiệp vụ (Business Requirements)}
\label{tab:3-br}
\small
\begin{tblr}{
  colspec = {X[0.9,l]X[2,l]X[5,l]X[1.2,l]}, rows = {abovesep=2pt, belowsep=2pt},
  hlines, vlines, row{1} = {font=\bfseries, bg=gray!20},
}
ID & Yêu cầu & Mô tả & Phụ trách \\
BR01 & Tập trung hóa giám sát nông nghiệp & Cung cấp góc nhìn hợp nhất về trang trại, thửa ruộng và thiết bị giám sát để quản lý nhất quán trên nhiều địa điểm. & Đạt \\
BR02 & Nâng cao nhận biết điều kiện đồng ruộng & Cung cấp kịp thời thông tin về môi trường và điều kiện đồng ruộng, giảm phụ thuộc vào đo thủ công. & \dien{...} \\
BR03 & Duy trì tính liên tục của giám sát & Giảm thiểu khoảng trống dữ liệu và bảo toàn thông tin khi thiết bị, kết nối hoặc dịch vụ gián đoạn tạm thời. & Nhân \\
BR04 & Thiết lập dữ liệu nông nghiệp đáng tin cậy & Bảo đảm thông tin nhất quán, tin cậy, truy vết được; giảm rủi ro do số đo không hợp lệ, trùng lặp hoặc gây hiểu nhầm. & \dien{...} \\
BR05 & Hỗ trợ quản lý trang trại dựa trên dữ liệu & Dùng số đo và xu hướng để định hướng tưới tiêu, phân bổ tài nguyên và quản lý thửa ruộng, giảm lãng phí. & Tịnh \\
BR06 & Đánh giá lịch sử và lập kế hoạch & Lưu trữ lịch sử để so sánh giữa các giai đoạn, nhận diện mùa vụ, đánh giá thực hành quản lý và lập kế hoạch. & \dien{...} \\
BR07 & Làm cho thông tin dễ tiếp cận & Cung cấp tóm tắt, xu hướng và so sánh thửa ruộng dễ hiểu cho quản lý và nhà phân tích. & Duy \\
BR08 & Hỗ trợ ứng phó sớm với rủi ro & Phát hiện sớm điều kiện bất thường và thiết bị ngừng hoạt động, theo dõi quá trình xử lý. & \dien{...} \\
\end{tblr}
\end{table}

\section{Nhu cầu sử dụng dữ liệu}
\label{sec:3-data-needs}
% TODO: ánh xạ BR -> loại phân tích -> độ trễ cần thiết

\section{Nguồn dữ liệu và đặc trưng dữ liệu}
\label{sec:3-data-sources}

\subsection{Nguồn dữ liệu}
% TODO: cảm biến độ ẩm đất, nhiệt độ, thời tiết, thiết bị tưới, danh mục trang trại/thửa ruộng

\subsection{Đặc trưng dữ liệu}
% TODO: tốc độ, khối lượng, dữ liệu đến muộn/trùng/thiếu, chất lượng

\section{Ánh xạ yêu cầu nghiệp vụ sang yêu cầu xử lý dữ liệu}
\label{sec:3-mapping}
% TODO
